\section{Representation Learning of Functional Brain Dynamics}

Representation learning has been applied to functional brain data to obtain lower-dimensional descriptions of complex brain activity. 
Sanz Perl et al. used variational autoencoders (VAEs) to create low-dimensional embeddings of collective brain dynamics~\cite{SanzPerl2020}. 
Their results showed that the learned latent space can preserve meaningful differences between brain states, supporting the use of generative models for studying the organization of high-dimensional brain dynamics.

Representation learning has also been applied directly to resting-state fMRI (rs-fMRI). 
Kim et al.\ used a convolutional VAE to represent individual fMRI activity patterns in a low-dimensional latent space~\cite{Kim2021}. 
As each fMRI frame is represented separately, changes in the latent representation over time form trajectories that describe the temporal evolution of brain activity. 
The learned representations were related to patterns of cortical activity and functional connectivity, showing that information about functional brain organization can be retained after dimensionality reduction.

Low-dimensional representations can also be used to compare different global brain states. 
Sanz Perl et al.\ applied deep generative models to fMRI data from different states of consciousness and found that these states showed an ordered organization within a low-dimensional space~\cite{SanzPerl2023}. 
This demonstrates that latent-space structure can reflect differences in functional brain dynamics between conditions. 
Together, these studies show that representation learning can provide compact descriptions of high-dimensional fMRI activity while preserving information about temporal dynamics and differences between brain states.
